% GENERATED FILE. Edit canonical lessons, then run npm run generate:books.
% canonical-source-hash: fnv1a64:8f0e4b6c1f8cafc5
% canonical-lessons: RU-C498-druzhit, RU-C498-druzhba, RU-C498-priyatel, RU-C498-rovesnik, RU-C498-odnoklassnik

\chapter{Friends and Classmates}
\label{ch:ru-friends-and-classmates}
\hlchaptermodality{498}

\begin{chapteropening}
\textbf{By the end of this chapter:} \emph{I can talk about friendship and people my own age.}

Complete the last lesson of chapter 498: I can talk about friendship and people my own age.
\end{chapteropening}

\section[\texorpdfstring{\ru{дружить} (druzhít)}{дружить (druzhít)}]{\ru{дружить} (druzhít) — to be friends}

\label{lesson:RU-C498-druzhit}



\begin{quote}
Before the new one: say the Russian for to be naughty, then the Russian for disobedient.
\end{quote}

\subsection*{You'll want to know: \ru{дружить}}
\textbf{\ru{дружить}} (\emph{druzhít}) — \textquotedblleft{}to be friends\textquotedblright{}.

\begin{quote}
\ru{Мы} \ru{дружим} \ru{с} \ru{детства}.
\end{quote}

\begin{quote}
\emph{My drúzhim s détstva.}
\end{quote}

\begin{quote}
We've been friends since childhood.
\end{quote}

Imperfective; present \ru{дружу}, \ru{дружишь}. \emph{\ru{Дружить} \ru{с} \ru{кем}}.

\begin{cousinweb}
\emph{\ru{дружить}} is built on \emph{\ru{друг}} 'friend', which is inherited from Proto-Slavic \emph{drug\ru{ъ}} 'friend' (EtymDB).
\end{cousinweb}

\begin{grammarlens}[title={What you've built}]
One more word for telling, explaining and describing at B1.
\end{grammarlens}

\subsection*{Your turn}
\begin{itemize}
\raggedright
  \item \emph{Say it:} \emph{druzhít}
  \item \emph{Say it:} \emph{druzhít}, once more
  \item \emph{Say it:} a sentence of your own with \emph{druzhít}
  \item \emph{Recall:} say the Russian for to be naughty, then the Russian for disobedient, then say \emph{druzhít} again
\end{itemize}

\begin{tcolorbox}[breakable,colback=teal!4,colframe=teal!35!black,title={Before you move on}]
What does \ru{дружить} mean? (\textquotedblleft{}To be friends\textquotedblright{}.) Say it once more.
\end{tcolorbox}
\section[\texorpdfstring{\ru{дружба} (drúzhba)}{дружба (drúzhba)}]{\ru{дружба} (drúzhba) — friendship}

\label{lesson:RU-C498-druzhba}



\begin{quote}
Before the new one: say the Russian for disobedient, then the Russian for to be friends.
\end{quote}

\subsection*{You'll want to know: \ru{дружба}}
\textbf{\ru{дружба}} (\emph{drúzhba}) — \textquotedblleft{}friendship\textquotedblright{}.

\begin{quote}
\ru{Наша} \ru{дружба} \ru{началась} \ru{в} \ru{школе}.
\end{quote}

\begin{quote}
\emph{Násha drúzhba nachalás v shkóle.}
\end{quote}

\begin{quote}
Our friendship began at school.
\end{quote}

Feminine.

\begin{cousinweb}
\emph{\ru{дружба}} is inherited from Proto-Slavic \emph{druž\ru{ь}ba}, built on \emph{drug\ru{ъ}} 'friend', which EtymDB traces to PIE \emph{d\textsuperscript{h}rewg\textsuperscript{h}-.}
\end{cousinweb}

\begin{grammarlens}[title={What you've built}]
One more word for telling, explaining and describing at B1.
\end{grammarlens}

\subsection*{Your turn}
\begin{itemize}
\raggedright
  \item \emph{Say it:} \emph{drúzhba}
  \item \emph{Say it:} \emph{drúzhba}, once more
  \item \emph{Say it:} a sentence of your own with \emph{drúzhba}
  \item \emph{Recall:} say the Russian for disobedient, then the Russian for to be friends, then say \emph{drúzhba} again
\end{itemize}

\begin{tcolorbox}[breakable,colback=teal!4,colframe=teal!35!black,title={Before you move on}]
What does \ru{дружба} mean? (\textquotedblleft{}Friendship\textquotedblright{}.) Say it once more.
\end{tcolorbox}
\section[\texorpdfstring{\ru{приятель} (priyátel)}{приятель (priyátel)}]{\ru{приятель} (priyátel) — a pal, a friend}

\label{lesson:RU-C498-priyatel}



\begin{quote}
Before the new one: say the Russian for to be friends, then the Russian for friendship.
\end{quote}

\subsection*{You'll want to know: \ru{приятель}}
\textbf{\ru{приятель}} (\emph{priyátel}) — \textquotedblleft{}a pal, a friend\textquotedblright{}.

\begin{quote}
\ru{Вчера} \ru{я} \ru{встретил} \ru{старого} \ru{приятеля}.
\end{quote}

\begin{quote}
\emph{Vcherá ya vstrétil stárogo priyátelya.}
\end{quote}

\begin{quote}
Yesterday I ran into an old friend.
\end{quote}

Masculine; feminine \ru{приятельница}. Less close than \emph{\ru{друг}}.

\begin{cousinweb}
\emph{\ru{приятель}} is inherited from Old East Slavic \emph{prijatel\ru{ь}}, from Proto-Slavic \emph{prijateľ\ru{ь}}, derived from \emph{prijati} 'to please, to favour' (EtymDB).
\end{cousinweb}

\begin{grammarlens}[title={What you've built}]
One more word for telling, explaining and describing at B1.
\end{grammarlens}

\subsection*{Your turn}
\begin{itemize}
\raggedright
  \item \emph{Say it:} \emph{priyátel}
  \item \emph{Say it:} \emph{priyátel}, once more
  \item \emph{Say it:} a sentence of your own with \emph{priyátel}
  \item \emph{Recall:} say the Russian for to be friends, then the Russian for friendship, then say \emph{priyátel} again
\end{itemize}

\begin{tcolorbox}[breakable,colback=teal!4,colframe=teal!35!black,title={Before you move on}]
What does \ru{приятель} mean? (\textquotedblleft{}A pal, a friend\textquotedblright{}.) Say it once more.
\end{tcolorbox}
\section[\texorpdfstring{\ru{ровесник} (rovésnik)}{ровесник (rovésnik)}]{\ru{ровесник} (rovésnik) — someone of the same age}

\label{lesson:RU-C498-rovesnik}



\begin{quote}
Before the new one: say the Russian for friendship, then the Russian for a pal.
\end{quote}

\subsection*{You'll want to know: \ru{ровесник}}
\textbf{\ru{ровесник}} (\emph{rovésnik}) — \textquotedblleft{}someone of the same age\textquotedblright{}.

\begin{quote}
\ru{Мы} \ru{с} \ru{ним} \ru{ровесники}.
\end{quote}

\begin{quote}
\emph{My s nim rovésniki.}
\end{quote}

\begin{quote}
He and I are the same age.
\end{quote}

Masculine; feminine \ru{ровесница}.

\begin{cousinweb}
\emph{\ru{ровесник}} was borrowed from Polish \emph{rówieśnik} 'a person of the same age' (EtymDB).
\end{cousinweb}

\begin{grammarlens}[title={What you've built}]
One more word for telling, explaining and describing at B1.
\end{grammarlens}

\subsection*{Your turn}
\begin{itemize}
\raggedright
  \item \emph{Say it:} \emph{rovésnik}
  \item \emph{Say it:} \emph{rovésnik}, once more
  \item \emph{Say it:} a sentence of your own with \emph{rovésnik}
  \item \emph{Recall:} say the Russian for friendship, then the Russian for a pal, then say \emph{rovésnik} again
\end{itemize}

\begin{tcolorbox}[breakable,colback=teal!4,colframe=teal!35!black,title={Before you move on}]
What does \ru{ровесник} mean? (\textquotedblleft{}Someone of the same age\textquotedblright{}.) Say it once more.
\end{tcolorbox}
\section[\texorpdfstring{\ru{одноклассник} (odnoklássnik)}{одноклассник (odnoklássnik)}]{\ru{одноклассник} (odnoklássnik) — a classmate}

\label{lesson:RU-C498-odnoklassnik}



\begin{quote}
Before the new one: say the Russian for a pal, then the Russian for someone of the same age.
\end{quote}

\subsection*{You'll want to know: \ru{одноклассник}}
\textbf{\ru{одноклассник}} (\emph{odnoklássnik}) — \textquotedblleft{}a classmate\textquotedblright{}.

\begin{quote}
\ru{В} \ru{субботу} \ru{я} \ru{ходил} \ru{на} \ru{встречу} \ru{с} \ru{одноклассниками}.
\end{quote}

\begin{quote}
\emph{V subbótu ya khodíl na vstréchu s odnoklássnikami.}
\end{quote}

\begin{quote}
On Saturday I went to a reunion with my classmates.
\end{quote}

Masculine; feminine \ru{одноклассница}. At university: \emph{\ru{однокурсник}}.

\begin{cousinweb}
\emph{\ru{одноклассник}} is a compound of \emph{\ru{один}} 'one' and \emph{\ru{класс}} 'class', which EtymDB traces to Latin \emph{classis}, with \emph{-\ru{ник}}; EtymDB has no entry for \emph{\ru{одноклассник}}.
\end{cousinweb}

\begin{grammarlens}[title={What you've built}]
One more word for telling, explaining and describing at B1.
\end{grammarlens}

\subsection*{Your turn}
\begin{itemize}
\raggedright
  \item \emph{Say it:} \emph{odnoklássnik}
  \item \emph{Say it:} \emph{odnoklássnik}, once more
  \item \emph{Say it:} a sentence of your own with \emph{odnoklássnik}
  \item \emph{Recall:} say the Russian for a pal, then the Russian for someone of the same age, then say \emph{odnoklássnik} again
\end{itemize}

\begin{tcolorbox}[breakable,colback=teal!4,colframe=teal!35!black,title={Before you move on}]
What does \ru{одноклассник} mean? (\textquotedblleft{}A classmate\textquotedblright{}.) Say it once more.
\end{tcolorbox}
